\documentclass{article}

\usepackage[utf8]{inputenc}
\usepackage[T1]{fontenc}
\usepackage{amsmath,amssymb,amsthm}
\usepackage{mathtools}
\usepackage{hyperref}
\usepackage{booktabs}

\makeatletter
\newcommand{\citep}[2][]{%
  \begingroup
  \def\@tmpkey{#2}%
  \ifx\@tmpkey\@empty\else
    (\ifx\@empty#1\@empty\else #1, \fi
     \@nameuse{bibkey@\@tmpkey})%
  \fi
  \endgroup
}
\expandafter\def\csname bibkey@lasser2026exo\endcsname{Exogenous~Verification}
\expandafter\def\csname bibkey@lasser2026phase\endcsname{Phase~Redundancy}
\expandafter\def\csname bibkey@lasser2026micro\endcsname{Microfoundation}
\expandafter\def\csname bibkey@lasser2026horizon\endcsname{Horizon~Aware}
\makeatother

\newtheorem{definition}{Definition}
\newtheorem{lemma}{Lemma}
\newtheorem{theorem}{Theorem}
\newtheorem{remark}{Remark}
\newtheorem{assumption}{Assumption}

\newcommand{\Aadv}{A_{\mathrm{adv}}}
\newcommand{\deltaadv}{\delta_{\mathrm{adv}}}
\newcommand{\Witness}{\mathcal{W}}
\newcommand{\Ledger}{\mathcal{L}}
\newcommand{\rext}{r_{\mathrm{ext}}}
\newcommand{\mindep}{m_{\mathrm{eff}}^{\mathrm{indep}}}
\newcommand{\deltastar}{\delta_{*}}
\newcommand{\Tbeta}{T_{\beta}}

\title{Lemma 5e (Environment-side witness extension): Full Proof Draft v1}
\author{Paper 10 working draft for codex review}
\date{}

\begin{document}

\maketitle

\section{Goal}

Promote the environment-side detection KL floor derivation
(currently an inline sketch in Paper~10 §4.11) to a full appendix
proof, mirroring the structure of Lemma~5b's appendix proof
(\texttt{appendices/proofs.tex} §A.5).

V1 of the inline statement asserts that the SPRT machinery applies
symmetrically to environment-side observables under $I_8$, yielding
a KL floor $\deltaadv^{\mathrm{env}} > 0$ over $V_{\mathrm{env}}$.
This draft formalizes both (a) the symmetric construction (mirror of
\citep[Exogenous~Verification]{lasser2026exo}'s agent-side
substrate-exclusivity) and (b) the per-variable KL floor derivations
for the four canonical environment-side observables.

\section{Setup: environment-side substrate-exclusive witnesses}

\subsection{Symmetric construction (mirror of Paper 5)}

\citep[Exogenous~Verification]{lasser2026exo} constructs
substrate-exclusive witnesses for agent-side observables: a witness
on substrate $s' \neq s(\mathrm{agent})$ attests to events on
substrate $s(\mathrm{agent})$ via cross-substrate verification.
The construction relies on three properties:
\begin{enumerate}
\item \emph{Substrate distinctness:} the witness substrate $s'$ is
  causally separable from the observed substrate $s$.
\item \emph{Cryptographic commitment:} witness records are committed
  via Pedersen commitments \emph{(or analogous binding scheme)}
  ensuring tamper-evidence.
\item \emph{Ledger publication:} commitments are published to the
  verification ledger $\Ledger$ under
  \citep[Exogenous~Verification]{lasser2026exo}'s
  substrate-exclusivity discipline.
\end{enumerate}

Invariant $I_8$ (Paper~10 §3) extends this construction to
environment-side observables.  The mirror property is:
\[
  I_8: \quad \forall v \in V_{\mathrm{env}}:\
  \exists\ \Witness_v \text{ on } s_{\mathrm{env}} \neq
  s(\mathrm{agent}).
\]
The witness substrate $s_{\mathrm{env}}$ is environment-side; the
witnessed quantity $v$ is an exogenous state variable; the
observed-substrate constraint becomes $s(\mathrm{agent})$.

\textbf{Construction note (new machinery).}  Paper~5 specifies the
agent-side substrate-exclusivity construction; Paper~10 §3 specifies
that environment-side substrate-exclusive witnesses use a separate
substrate partition with separate trust models, but inherits the
cryptographic and ledger-publication discipline from Paper~5.
Substantively new machinery: (i) the environment-side substrate
partition, (ii) trusted-setup ceremonies for environment-side
witnesses, (iii) coverage specification by enumerating
$V_{\mathrm{env}}$ in deployment documentation.

\subsection{The detection class $\Aadv^{\mathrm{env}}$}

\begin{definition}[Environment-side adversarial class
$\Aadv^{\mathrm{env}}$]
A regime-(iii) strategy $q$ achieving its effect via environment
manipulation belongs to $\Aadv^{\mathrm{env}}$ if and only if $q$
produces an observable shift in at least one $v \in V_{\mathrm{env}}$
under the SPRT alternative versus baseline.
\end{definition}

\textbf{The named residual.}  Environment manipulations targeting
exogenous variables outside $V_{\mathrm{env}}$ (variables not
monitored by environment-side witnesses) are outside
$\Aadv^{\mathrm{env}}$ and are not detected by Lemma~5e.  This is
Theorem~1 residual~(R2).  The residual is real: $V_{\mathrm{env}}$
is finite by deployment specification, and adversarial environment
manipulation that perturbs only unmonitored exogenous variables
escapes detection.

\section{Per-variable KL floor derivations}

We derive $\deltaadv^{\mathrm{env}}$ as the minimum over four
per-variable KL floors, one for each canonical environment-side
observable enumerated in Paper~10's deployment-tooling
specification (\S8.1) and $I_8$'s threshold semantics:
\begin{itemize}
\item $V_1$: cross-substrate cooperative novelty rate $\rext$
  (Poisson).
\item $V_2$: substrate-distinctness independence indicator
  (Bernoulli).
\item $V_3$: adversarial-event arrival rate $\lambda$ (Poisson).
\item $V_4$: trusted-setup status flag (Bernoulli).
\end{itemize}

For each variable, we (a) state the baseline distribution
$p_0^{\mathrm{env}}(v)$, (b) the least-favorable adversarial
alternative $p_1^{\mathrm{env}}(v)$ subject to the SPRT-detectability
constraint, and (c) the resulting KL divergence floor.

\subsection{Variable $V_1$: cross-substrate cooperative rate
(Poisson)}

\textbf{Baseline.}  Under stance S0,
\citep[Horizon~Aware]{lasser2026horizon}'s cross-substrate
cooperative novelty rate $\rext > 0$ is environment-observable:
ledger entries record cross-substrate cooperative events with
Poisson arrival rate $\lambda_0^{(1)}$ in the no-adversary baseline.
The witness $\Witness_{V_1}$ on $s_{\mathrm{env}}$ records each
arrival with timestamp and substrate-pair attestation.

\textbf{Adversarial alternative.}  An environment-manipulating
strategy that suppresses cross-substrate cooperatives (e.g., by
disrupting environmental conditions for cross-substrate cooperation)
reduces the rate to at most $\lambda_0^{(1)} - \eta_1$ for some
$\eta_1 \in (0, \lambda_0^{(1)})$, where $\eta_1$ is the
SPRT-detectability threshold.

\textbf{KL divergence.}  The Poisson KL between rate-$\lambda_1$ and
rate-$\lambda_0$ processes per unit time is:
\[
  D_{\mathrm{KL}}(\mathrm{Poi}(\lambda_1) \,\|\,
  \mathrm{Poi}(\lambda_0)) \;=\; \lambda_1 \log(\lambda_1 /
  \lambda_0) - \lambda_1 + \lambda_0.
\]
This is convex in $\lambda_1$ and increasing in
$|\lambda_1 - \lambda_0|$.  The least-favorable choice within the
constraint $\lambda_1 \leq \lambda_0^{(1)} - \eta_1$ is
$\lambda_1 = \lambda_0^{(1)} - \eta_1$, giving:
\[
  \delta_1^{\mathrm{env}} \;=\; (\lambda_0^{(1)} - \eta_1)
  \log\!\frac{\lambda_0^{(1)} - \eta_1}{\lambda_0^{(1)}} -
  (\lambda_0^{(1)} - \eta_1) + \lambda_0^{(1)} \;>\; 0.
\]

\subsection{Variable $V_2$: substrate-distinctness independence
indicator (Bernoulli)}

\textbf{Baseline.}  Per-period substrate-distinctness independence
indicator $\xi_t \in \{0, 1\}$ is recorded by $\Witness_{V_2}$:
$\xi_t = 1$ if the period-$t$ failure-correlation audit verifies
$\mindep \geq m^*$ (independence holds), $\xi_t = 0$ otherwise.
Baseline rate: $\Pr[\xi_t = 1] = p_0^{(2)} \in (0, 1]$ (typically
near $1$ for well-anchored deployments).

\textbf{Adversarial alternative.}  An environment-manipulating
strategy that erodes substrate distinctness (e.g., by introducing
correlated failure modes across nominally distinct substrates)
reduces the success rate to at most $p_0^{(2)} - \eta_2$ for some
$\eta_2 \in (0, p_0^{(2)}]$.

\textbf{KL divergence.}  Bernoulli KL:
\[
  \delta_2^{\mathrm{env}} \;=\; (p_0^{(2)} - \eta_2)
  \log\!\frac{p_0^{(2)} - \eta_2}{p_0^{(2)}} +
  (1 - p_0^{(2)} + \eta_2) \log\!\frac{1 - p_0^{(2)} + \eta_2}
  {1 - p_0^{(2)}} \;>\; 0.
\]
This is the least-favorable choice within the constraint
$p_1 \leq p_0^{(2)} - \eta_2$.

\subsection{Variable $V_3$: adversarial-event arrival rate
(Poisson)}

\textbf{Baseline.}  Adversarial events arrive on the environment
side with Poisson rate $\lambda_0^{(3)} \geq 0$ (the baseline
threat-model arrival rate, calibrated by the deployment's threat
model).  Witness $\Witness_{V_3}$ records arrivals.

\textbf{Adversarial alternative.}  An adversarial strategy that
\emph{increases} the arrival rate to $\lambda_0^{(3)} + \eta_3$
(introducing additional adversarial events) is detectable if
$\eta_3 > 0$ exceeds the SPRT detection threshold.

\textbf{KL divergence.}  Same Poisson form as $V_1$, but with
positive shift $\eta_3$:
\[
  \delta_3^{\mathrm{env}} \;=\; (\lambda_0^{(3)} + \eta_3)
  \log\!\frac{\lambda_0^{(3)} + \eta_3}{\lambda_0^{(3)}} -
  (\lambda_0^{(3)} + \eta_3) + \lambda_0^{(3)} \;>\; 0.
\]
The case $\lambda_0^{(3)} = 0$ (no baseline adversarial events)
requires special handling: the KL is then $\eta_3 \log(\eta_3 / 0)
= \infty$ formally, but operationally is bounded by the
environment-side substrate's measurement resolution.  Deployments
should use $\lambda_0^{(3)} > 0$ as a regularized baseline.

\subsection{Variable $V_4$: trusted-setup status flag (Bernoulli)}

\textbf{Baseline.}  The trusted-setup status flag $\sigma_t \in
\{0, 1\}$ records whether the cryptographic trusted-setup ceremony
remains intact at period $t$: $\sigma_t = 1$ if integrity audits
pass, $\sigma_t = 0$ if compromise is detected.  Baseline rate
$p_0^{(4)} \in (0, 1]$ (near $1$ for properly maintained setups).

\textbf{Adversarial alternative.}  Environment-side adversarial
strategies that compromise the trusted-setup substrate reduce the
flag's success rate to $p_0^{(4)} - \eta_4$.

\textbf{KL divergence.}  Same Bernoulli form as $V_2$:
\[
  \delta_4^{\mathrm{env}} \;=\; (p_0^{(4)} - \eta_4)
  \log\!\frac{p_0^{(4)} - \eta_4}{p_0^{(4)}} +
  (1 - p_0^{(4)} + \eta_4) \log\!\frac{1 - p_0^{(4)} + \eta_4}
  {1 - p_0^{(4)}} \;>\; 0.
\]

\subsection{Combined floor}

\[
  \deltaadv^{\mathrm{env}} \;:=\; \min\{\delta_1^{\mathrm{env}},
  \delta_2^{\mathrm{env}}, \delta_3^{\mathrm{env}},
  \delta_4^{\mathrm{env}}\} \;>\; 0.
\]
Each $\delta_v^{\mathrm{env}} > 0$ since each least-favorable
adversarial shift is strictly different from the baseline; the
minimum is therefore strictly positive.

\section{Statement of Lemma 5e (full)}

\begin{lemma}[Environment-side witness extension]
\label{lem:5e}
Under invariant $I_8$ (environment-side substrate-exclusive
witnesses on exogenous variables), Lemma~5b's channel-restricted
detection class extends to environment-side observables.
Specifically, for any $q \in \Aadv^{\mathrm{env}}$:
\[
  D_{\mathrm{KL}}(q \,\|\, p_0^{\mathrm{env}}) \;\geq\;
  \deltaadv^{\mathrm{env}} \;>\; 0,
\]
where $\deltaadv^{\mathrm{env}} = \min\{\delta_v^{\mathrm{env}} :
v \in V_{\mathrm{env}}\}$ is the minimum over per-variable KL
floors derived from environment-side least-favorable alternatives.

For the canonical four-variable $V_{\mathrm{env}} = \{V_1, V_2,
V_3, V_4\}$ enumerated in §3, the explicit per-variable floors are:
\begin{itemize}
\item $\delta_1^{\mathrm{env}}$ (Poisson cross-substrate cooperative
  rate),
\item $\delta_2^{\mathrm{env}}$ (Bernoulli substrate-distinctness
  indicator),
\item $\delta_3^{\mathrm{env}}$ (Poisson adversarial-event arrival
  rate),
\item $\delta_4^{\mathrm{env}}$ (Bernoulli trusted-setup status
  flag),
\end{itemize}
each $> 0$ as derived in §3.1--3.4.
\end{lemma}

\begin{proof}
The proof has three parts: (a) the symmetric construction (§2),
(b) the per-variable KL floor derivations (§3), and (c) the
combined floor.

(a)~The symmetric construction inherits Paper~5's
substrate-exclusivity machinery (cryptographic commitment +
ledger publication + cross-substrate verification) with the
substrate-roles swapped: $s_{\mathrm{env}}$ replaces $s'$
(witness substrate), and exogenous variables $v$ replace agent
observables.  $I_8$ requires this construction for every $v \in
V_{\mathrm{env}}$.

(b)~For each $v_i \in V_{\mathrm{env}}$, the SPRT machinery applies
on the witness-recorded observable.  The least-favorable
adversarial alternative within the SPRT-detectability constraint
gives a strictly positive per-variable KL floor
$\delta_i^{\mathrm{env}} > 0$ (§3.1--3.4 derivations).

(c)~$\deltaadv^{\mathrm{env}} = \min_i \delta_i^{\mathrm{env}}$ is
strictly positive as the minimum of finitely many strictly
positive quantities.  $\Box$
\end{proof}

\section{Discussion}

\paragraph{Composition with Lemma 6.}  The post-clipping union-class
drift floor $\deltastar$ used in Lemma~6's $\Tbeta$ derivation is
$\deltastar \leq \min(\deltaadv, \deltaadv^{\mathrm{env}})$, with
the inequality accounting for any drift reduction from (C5.HOEFF)
clipping.  Lemma~5e supplies $\deltaadv^{\mathrm{env}}$ for the
union; Lemma~5b supplies $\deltaadv$.

\paragraph{Coverage and the named residual (R2).}
$V_{\mathrm{env}}$ is finite by deployment specification.  Variables
outside $V_{\mathrm{env}}$ are real exogenous quantities not
monitored by environment-side witnesses; environment manipulations
targeting only those unmonitored variables fall into Theorem~1
residual~(R2) and are outside the deployment-safety guarantee.
The deployment must enumerate $V_{\mathrm{env}}$ exhaustively per
its threat model; under-coverage shifts adversarial mass into
(R2).

\paragraph{Asymmetric vs. symmetric machinery.}
Paper~5's agent-side substrate-exclusivity is constructed; Paper~10
§3 specifies that environment-side substrate-exclusivity is the
mirror, but is new machinery (separate substrate partition, separate
trust models).  This lemma's symmetric-construction argument is
\emph{structural}: it inherits Paper~5's discipline.  The
deployment must instantiate the environment-side construction
(Audit~2/3 in §8 deployment tooling).

\paragraph{Choice of $\eta_v$.}
Each per-variable shift threshold $\eta_v$ is the SPRT
detectability threshold for that variable.  Operationally, $\eta_v$
is set by the deployment's threat model: how much shift in $v$
must be detected for the deployment to consider the manipulation
adversarially material.  Tighter $\eta_v$ gives smaller
$\delta_v^{\mathrm{env}}$ (higher detection sensitivity but lower
KL floor); looser $\eta_v$ gives larger $\delta_v^{\mathrm{env}}$
(lower sensitivity but higher floor, hence faster detection in
$\Tbeta$).

\section{What this proof does and does not establish}

\textbf{Establishes:}
\begin{itemize}
\item The symmetric-construction argument formalizing $I_8$'s
  extension of Paper~5's substrate-exclusivity to environment-side
  observables.
\item Per-variable KL floor derivations for the four canonical
  environment-side observables (Poisson for rates; Bernoulli for
  binary flags), each strictly positive.
\item The combined floor $\deltaadv^{\mathrm{env}} = \min_v
  \delta_v^{\mathrm{env}} > 0$ as a deployment-class constant.
\item Composition with Lemma~6 via the post-clipping union-class
  drift floor $\deltastar$.
\end{itemize}

\textbf{Does not establish (deferred):}
\begin{itemize}
\item Tightness of per-variable shifts $\eta_v$.  These are
  threat-model-derived and deployment-policy-set.
\item Exhaustive coverage of $V_{\mathrm{env}}$.  This is a
  deployment-policy enumeration question; under-coverage falls into
  residual~(R2).
\item Security proof of the environment-side substrate-exclusivity
  construction.  Paper~10 §3 specifies the construction, not the
  security proof; cryptographic security is inherited from
  Paper~5's substrate-exclusivity discipline.
\item Multinomial extension.  $V_{\mathrm{env}}$ is enumerated as
  Poisson + Bernoulli observables for the canonical four; if
  deployment policy adds multinomial environment-side observables
  (e.g., partition over trade-flow categories), the analogous
  multinomial KL derivation from Lemma~5b's appendix proof applies.
\end{itemize}

\section{Specific verification questions for codex review}

\begin{enumerate}
\item Is the symmetric-construction argument rigorous, or does the
  ``mirror of Paper~5'' framing leave operational gaps (e.g.,
  ledger-commitment trust assumptions for environment-side
  witnesses that don't transfer cleanly from agent-side)?

\item Is the per-variable KL floor derivation complete?  In
  particular: $V_3$ has the case $\lambda_0^{(3)} = 0$ (no
  baseline adversarial events) where the Poisson KL formally
  diverges; the draft addresses this with a regularization note,
  but is the regularization sufficient or does it warrant a
  separate clause?

\item Is $V_{\mathrm{env}} = \{V_1, V_2, V_3, V_4\}$ canonical?
  The deployment-tooling enumeration mentions four named
  observables (cooperative rate, distinctness indicator,
  adversarial arrival rate, trusted-setup flag), but should the
  proof allow for additional deployment-specific variables?  If so,
  the per-variable derivation pattern (Poisson/Bernoulli/
  multinomial) covers them.

\item Is the named residual (R2) connection correct?  The proof
  asserts that variables outside $V_{\mathrm{env}}$ are residuals;
  is this consistent with Theorem~1's residual~(R2) statement, or
  does the residual definition need refining?

\item Does the proof need Theorem-level conditions analogous to
  (C5.SPRT) bundle?  Specifically, the SPRT-applicability
  conditions ((C5.HOEFF) clipping, (C5.MULT) cardinality, (C5.IID)
  adapted increments, (C5.SUPP) finiteness) apply to both
  agent-side and environment-side LLR processes; should the
  conditions be restated as covering both, or is the
  agent-vs-environment-side symmetry implicit?
\end{enumerate}

\section{Structural questions for the v-cycle}

Following the workflow used for Lemmas~1, 2, 5a, 6:

\textbf{Q-A.} Is this proof appendix-worthy, or does it stay
inline?  TODO\_proofs.md says ``Body for the symmetry argument;
appendix for the KL floor derivation.''  V1 has both in one
draft for review.

\textbf{Q-B.} What's the right page-budget?  Lemma~5b's appendix
proof is $\sim$1.5pp; mirroring that, Lemma~5e's appendix proof
should be similar.  V1 is more detailed than necessary; the final
version should be tight.

\textbf{Q-C.} Are there theorem-level conditions to surface?
Lemmas~1, 2, 5a, 6 each surfaced new conditions (C6, C7, C8, C9,
C10, C11, C5.SPRT bundle, C11.CLK).  Lemma~5e may or may not
require new conditions; the v-cycle should determine this.

\end{document}
